% ═══════════════════════════════════════════════════════════════
% LEHRERHINWEISE: Strahlenschutz (Physik Klasse 10)
% Stunde 9-10 (90 Minuten)
% ═══════════════════════════════════════════════════════════════

\documentclass[11pt, a4paper]{article}

\usepackage[utf8]{inputenc}
\usepackage[T1]{fontenc}
\usepackage[ngerman]{babel}
\usepackage{helvet}
\renewcommand{\familydefault}{\sfdefault}

\usepackage[margin=2cm]{geometry}
\usepackage{enumitem}
\usepackage{tabularx}
\usepackage{booktabs}
\usepackage{array}
\usepackage{xcolor}
\usepackage{tcolorbox}
\usepackage{fancyhdr}

\setlist{nosep, leftmargin=1.5em}

% Farben
\definecolor{physik}{HTML}{2c5aa0}
\definecolor{grau}{HTML}{888888}
\definecolor{hellgrau}{HTML}{f0f0f0}
\definecolor{warnung}{HTML}{b35c00}
\definecolor{warnungbg}{HTML}{fff3e0}

% Kopfzeile
\pagestyle{fancy}
\fancyhf{}
\renewcommand{\headrulewidth}{0.4pt}
\lhead{\small\textbf{LEHRERHINWEISE} -- Physik Klasse 10}
\rhead{\small Strahlenschutz (Stunde 9-10)}
\cfoot{\small\thepage}

% Boxen
\newtcolorbox{warnbox}[1][]{
    colback=warnungbg, colframe=warnung, boxrule=0.8pt, arc=0pt,
    left=6pt, right=6pt, top=6pt, bottom=6pt,
    fonttitle=\bfseries\small, title=#1
}

\newcommand{\abschnitt}[1]{%
    \vspace{10pt}\noindent\textbf{\textcolor{physik}{\large #1}}\vspace{6pt}\hrule\vspace{8pt}
}

\begin{document}

\begin{center}
{\LARGE\bfseries Lehrerhinweise: Strahlenschutz}\\[6pt]
{\large Physik Klasse 10 -- Stunde 9-10 (90 Min)}\\[4pt]
{\small\textcolor{grau}{Kernphysik-Einheit, Teil 5 (Abschluss)}}
\end{center}
\vspace{8pt}\hrule\vspace{12pt}

% ═══════════════════════════════════════════════════════════════
\abschnitt{1. Stundenverlauf (90 Minuten)}

\begin{tabularx}{\textwidth}{|c|c|X|l|}
\hline
\textbf{Zeit} & \textbf{Phase} & \textbf{Inhalt} & \textbf{Material} \\
\hline
0-5 & Einstieg & Frage: Was bedeutet das Strahlenwarnzeichen? & Bild \\
\hline
5-15 & Input & Lernpfad: Dosisbegriffe (Bq, Gy, Sv) & LP, Beamer \\
\hline
15-25 & Erarbeitung & AB Kasten 1-2: Einheiten, Grenzwerte & AB-05 \\
\hline
25-35 & Input & Lernpfad: Die 3 A's des Strahlenschutzes & LP \\
\hline
35-45 & Erarbeitung & AB Kasten 3, Übung 1-2 & AB-05 \\
\hline
45-60 & Fallstudien & AB Seite 2: Röntgenabteilung, Transport & AB-05 \\
\hline
60-70 & Berechnung & Abstandsgesetz (1/r²), Tabelle ausfüllen & AB-05 \\
\hline
70-80 & Diskussion & Übung 3-4: Nutzen vs. Risiko (Röntgen) & Plenum \\
\hline
80-90 & Sicherung & Zusammenfassung Gesamteinheit, Ausblick Test & Tafel \\
\hline
\end{tabularx}

\vspace{12pt}

% ═══════════════════════════════════════════════════════════════
\abschnitt{2. Differenzierung}

\begin{tabularx}{\textwidth}{|c|X|X|}
\hline
\textbf{Niveau} & \textbf{Fokus} & \textbf{Hilfestellung} \\
\hline
★ (BR) & 3 A's merken, Einheiten zuordnen & Merkblatt 3 A's \\
\hline
★★ (MR) & + Berechnungen mit Abstandsgesetz & Formelkarte, Taschenrechner \\
\hline
★★★ (GY) & + Fallstudien vollständig, Bewertung & Selbstständig \\
\hline
\end{tabularx}

\vspace{12pt}

% ═══════════════════════════════════════════════════════════════
\abschnitt{3. Häufige Fehler \& Reaktionen}

\begin{warnbox}[Typische Schülerfehler]
\begin{enumerate}
\item \textbf{Verwechslung Bq, Gy, Sv}\\
→ Merkhilfe: \textbf{B}q = \textbf{B}equerel = „wie viel strahlt?" (Aktivität). \textbf{S}v = \textbf{S}chaden (biol. Wirkung).

\item \textbf{„Jede Strahlung ist gefährlich, Röntgen ablehnen!"}\\
→ Diskussion: Nutzen-Risiko-Abwägung. Diagnose kann Leben retten! ALARA-Prinzip (As Low As Reasonably Achievable).

\item \textbf{Abstandsgesetz falsch angewendet}\\
→ Klärung: Doppelter Abstand = 1/4 der Dosis (nicht 1/2!). Es ist $1/r^2$, nicht $1/r$.

\item \textbf{„Natürliche Strahlung ist harmlos"}\\
→ Differenzieren: Radon in Kellern kann problematisch sein. Höhenstrahlung bei Flügen. Dosis ist Dosis, egal ob natürlich oder künstlich.
\end{enumerate}
\end{warnbox}

\vspace{12pt}

% ═══════════════════════════════════════════════════════════════
\abschnitt{4. Benötigtes Material}

\begin{itemize}
\item Lernpfad: \texttt{lernpfad-strahlenschutz.html}
\item Arbeitsblatt: \texttt{AB-05-strahlenschutz.pdf} (Klassensatz)
\item Musterlösung: \texttt{ML-05-strahlenschutz.pdf}
\item Bild: Strahlenwarnzeichen (Trefoil)
\item Optional: Dosimeter zum Zeigen
\item Optional: Video „Strahlenschutz im Krankenhaus"
\end{itemize}

\vspace{12pt}

% ═══════════════════════════════════════════════════════════════
\abschnitt{5. Lösungen (Kurzfassung)}

\textbf{Kasten 1:} Aktivität = Bq (Zerfälle/s), Energiedosis = Gy (J/kg), Äquivalentdosis = Sv (biol. Wirkung). 1 Sv = 1000 mSv = 1.000.000 µSv

\textbf{Kasten 2:} Natürlich: 2-3 mSv/Jahr. Grenzwert Bevölkerung: 1 mSv. Beruflich: 20 mSv. Röntgen-Thorax: 0,01-0,1 mSv. CT: 2-20 mSv. Strahlenkrankheit ab ca. 1 Sv.

\textbf{Kasten 3:} 3 A's: Abstand (größer = weniger), Abschirmung (Material dazwischen), Aufenthaltszeit (kürzer = weniger). Dosis $\sim 1/r^2$

\textbf{Übung 2a:} 4m = 2×2m → Faktor 4 weniger → 100/4 = 25 µSv/h

\textbf{Übung 2b:} 20 mSv ÷ 0,1 mSv/h = 200 Stunden

\textbf{Fall A:} 20×10µSv×250d = 50.000 µSv = 50 mSv > Grenzwert 20 mSv!

\textbf{Fall C:} 2m: 1/4 → 125 µSv/h / 5m: 1/25 → 20 µSv/h / 10m: 1/100 → 5 µSv/h

\vspace{8pt}
\textit{→ Vollständige Lösungen siehe ML-05-strahlenschutz.pdf}

\vspace{12pt}

% ═══════════════════════════════════════════════════════════════
\abschnitt{6. Hinweis: Abschlusstest}

Nach dieser Stunde ist die Kernphysik-Einheit abgeschlossen. In der Folgestunde kann der Test geschrieben werden:

\begin{itemize}
\item Online: \texttt{test-kernphysik-online.html}
\item Papier: \texttt{test-kernphysik-papier.pdf}
\end{itemize}

\textbf{Wiederholungsempfehlung:} SuS sollen alle AB-Kästen durchgehen und die Übungen nochmal rechnen.

\end{document}
